\section{Discussion}
\label{sec:discussion}

\subsection{Comparison with Galactic WC stars}
\label{sec:comparison_mw}

\citet{Dsilva2023} measured the observed spectroscopic binary fraction of
Galactic WC/WO stars using VLT/X-shooter spectroscopy and the same
cross-correlation framework adopted here, providing the natural
methodological twin of the present LMC measurement. Both surveys derive
radial velocities from the \primaryline emission line, classify a star as
a binary when $\Delta\mathrm{RV} > \drvthresh$\,\kms is recorded at
$\sigthresh\sigma$ significance, and apply Monte-Carlo bias correction to
recover the intrinsic binary fraction (Sect.~\ref{sec:bias}). The
methodological alignment removes the principal systematic differences
between published WR binary fractions, so the residual offset can be
attributed to the underlying stellar populations rather than to detection
choices. The cadence-dependent bias differs between the two samples: the
LMC targets were observed across \Nepochs epochs per star, whereas the
Galactic survey adopted per-target cadences tailored to expected orbital
periods. The corrected fraction \fbincorr, rather than \fbinobs (or
\fbinFinal), therefore provides the apples-to-apples comparison: a value
of \fbincorr above the Galactic estimate of \citet{Dsilva2023} would
indicate a stronger binary contribution at LMC metallicity, while a
comparable or lower value would constrain that prediction.

\begin{figure}
\centering
% \includegraphics[width=\columnwidth]{plots/lmc_vs_mw.pdf}
\caption{Comparison of the bias-corrected binary fraction and
  orbital-period distribution exponent for the LMC WC/WO sample (this
  work) and the Galactic WC/WO sample of \citet{Dsilva2023}. Error bars
  are 68\% HDIs from the marginalised log-likelihood. The horizontal
  dashed lines mark the Galactic reference values.}
\label{fig:lmc_vs_mw}
\end{figure}

\subsection{Metallicity dependence}
\label{sec:metallicity}

At the sub-solar metallicity of the LMC ($Z \approx 0.3$--$0.5\,\Zsun$;
\citealt{Russell1992}), line-driven mass loss is weaker than at solar
metallicity, raising the minimum initial mass at which a single star can
shed its hydrogen envelope and reach the WR stage \citep{Shenar2020}. A
generic prediction follows: as the single-star channel becomes
progressively inefficient, the binary channel must contribute a larger
fraction of the WC/WO population at LMC metallicity than at Milky Way
metallicity. Operationally, this expectation should manifest as
$f_\mathrm{bin,LMC} > f_\mathrm{bin,MW}$ once both samples are corrected
for observational bias. The relevant test is therefore a comparison of
the corrected LMC fraction \fbincorr with the corrected Galactic fraction
of \citet{Dsilva2023}, both of which marginalise over the assumed
orbital-parameter distributions. The corrected \fbincorr provides the
test: a value above the corresponding Galactic estimate would support
the metallicity-dependent binary-channel prediction, while a value
below or comparable to the Galactic estimate would constrain it.
A second observable carries complementary information about formation
channels: the orbital-period exponent \pibestfit recovered from the
bias-correction grid (Sect.~\ref{sec:bias}). A bottom-heavy period
distribution (more long-period systems) is expected when wide-orbit mass
transfer dominates the production of stripped helium stars, whereas a
top-heavy distribution (more short-period systems) points to
Roche-lobe-overflow stripping in close binaries. The two observables
together --- $f_\mathrm{bin}$ and the period-distribution shape --- are
therefore the primary handles on the metallicity dependence accessible
from this work.

\subsection{Formation channels}
\label{sec:formation}

Two competing channels are typically invoked to produce WC/WO stars: (a)
single-star evolution in which strong line-driven winds strip the
hydrogen and helium envelopes of an initially very massive star
\citep[the Conti scenario;][]{Conti1976}, and (b) binary mass transfer in
which Roche-lobe overflow strips the envelope onto a companion, exposing
the carbon-rich core at lower initial masses than the single-star
pathway requires \citep{Paczynski1967}.
Population-synthesis calculations predict that the binary stripping
channel becomes increasingly dominant as metallicity decreases, because
the single-star channel is throttled by the metallicity dependence of
line-driven mass loss \citep{Marchant2023}. The corrected binary fraction
\fbincorr is the principal observable that discriminates between the two
channels. A value $\gtrsim 50\%$ favours a binary-dominated origin for
the LMC WC/WO population, a value $\lesssim 30\%$ favours single-star
evolution, and intermediate values imply that both channels contribute
non-negligibly. The orbital-period exponent \pibestfit, together with
its Langer-model alternative (Sect.~\ref{sec:langer_model}), provides an
orthogonal test by tracing the imprint of the dominant interaction mode
on the surviving orbital architecture.
The mass-ratio distribution, which is in principle a third diagnostic,
is not constrained by this work: the survey is RV-only with limited
orbital-phase coverage per star (Sect.~\ref{sec:threshold_results}), so
the discussion here stops at the binary fraction and the
period-distribution shape.

\subsection{WC + compact-object candidates}
\label{sec:compact_objects}

As a follow-up to the RV-based binary classification, \NcandidateBH
stars in the sample have been flagged as candidate WC + compact-object
(WC+BH) systems and are being targeted for phase-resolved spectroscopic
monitoring (proposal in preparation; results will be appended to a
future companion paper). The initial selection was guided by the
orbital phase relationship between absorption and emission lines: in a
genuine WR+OB binary, photospheric absorption from the OB companion is
expected to phase-lock with opposing radial-velocity signatures
relative to the WR wind emission lines, while WR + compact-object
systems show either no or anomalous absorption-line phase behaviour.
The final candidate selection further refined this signature using a
complementary diagnostic that will be described in the companion paper.
WC+BH systems at LMC metallicity directly probe the black-hole mass
distribution from stripped helium-star progenitors, complementing
population-synthesis predictions of compact-object merger rates
\citep{vandenHeuvel2017}.
